% Estrutura baseada em Wohlin et al. (escopo e planejamento do experimento).
\chapter[Planejamento do Experimento]{Planejamento do Experimento}
\label{cap:planejamento}

\section{Definição do escopo do experimento}
\label{sec:escopo}

\subsection{Objetivo}

De acordo com o levantamento bibliográfico efetuado, o RAG é bastante usado para elaborar histórias de usuário e critérios de aceitação \cite{rawson2025, rawson2025framework, rotar2025}, mas são escassas as pesquisas que confrontam a geração com e sem RAG \cite{wang2025}.

Sendo assim, o objetivo deste trabalho é testar de forma empírica se existe diferença estatisticamente significativa entre a acurácia dos critérios de aceitação em formato Gherkin gerados com e sem RAG.

\subsection{Questão de pesquisa}
\label{sec:quest-de-pesq}
Para obter uma resposta à questão de pesquisa principal, que foi apresentada na Seção~\ref{sec:questao-pesquisa} (\textit{Como os grandes modelos de linguagem têm contribuído na geração
automática de critérios de aceitação em formato Gherkin a partir da
documentação de histórias de usuário, e com que nível de acurácia?}), foi utilizada a metodologia GQM \cite{basili1994}, que permitiu dividir a questão de pesquisa principal em questões de pesquisa mais específicas e, para cada uma delas, definir métricas que podem ajudar na avaliação.

\begin{itemize}
    \item \textbf{Questão Específica 1}: 
    Existe diferença entre a acurácia dos critérios de aceitação em formato Gherkin gerados com \textit{RAG} e aqueles gerados sem RAG?

    Para medir a acurácia, os critérios de aceitação do gabarito serão divididos em unidades menores, chamadas \textit{objetivos de teste atômicos}, cada uma delas com o seu próprio resultado que pode ser verificado.
    
    \begin{itemize}
        \item \textbf{Métrica 1.1 (Acc@Hit):} Percentual de \textit{objetivos de teste atômicos} de cada história de usuário que são parcialmente contemplados pelos critérios de aceitação gerados.  
        \item \textbf{Métrica 1.2 (Acc@Cor):} Percentual de \textit{objetivos de teste atômicos} de cada história de usuário que são totalmente contemplados pelos critérios de aceitação gerados.
    \end{itemize}
\end{itemize}

\subsection{Trabalhos relacionados}
% {aqui vão os estudos da revisão mais próximos do nosso e a diferença deles para o nosso, porque posiciona a contribuição}
%\cite{wohlin2000}
%\cite{wang2025}
%fazer por último

Este trabalho se fundamenta em dois tipos de estudos: as referências metodológicas, que apresentaram as orientações para o planejamento do experimento, e o estudo de \citeonline{wang2025}, que forneceu a base de referência \textit{ground truth} para a definição dos tratamentos e da validação.

A referência metodológica utilizada foi o trabalho de \citeonline{wohlin2000}, que propõe uma estrutura para o processo experimental na Engenharia de Software, e que serviu como guia para a escolha do delineamento experimental, a formulação de hipóteses, a escolha de testes estatísticos e a discussão de ameaças à validade. Para a definição da pergunta de pesquisa específica e suas métricas a partir do objetivo proposto, utilizou-se a metodologia GQM \cite{basili1994}.

O trabalho de \citeonline{wang2025} propôs a arquitetura RAGcceptance M2RE para a elaboração de critérios de aceitação, que integra a recuperação de dados textuais e visuais, e que foi validada através de um estudo de ablação sem a utilização de RAG. Esse trabalho serviu como base para a adaptação do  \textit{pipeline} de recuperação e geração, o modelo de \textit{prompt}, a decomposição da base de referência \textit{ground truth} em objetivos atômicos e o comitê de LLMs avaliadores.

Diferentemente daquele estudo, em que o contraste entre o uso ou não de RAG faz parte da validação da proposta, aqui o RAG é o único fator experimental. O modelo, \textit{prompt}, parâmetros de geração e base de conhecimento são fixos, e o contraste entre os tratamentos é testado via um teste de hipóteses para amostras pareadas.

Portanto, esses estudos serviram como base metodológica e de ponto de partida para a elaboração dos tratamentos.

\section{Planejamento}
\label{sec:planejamento}

\subsection{Contexto}
% {aqui vão os projetos open source web ágeis escolhidos, os critérios de seleção e as histórias de usuário usadas, porque o contexto do GQM exige esse cenário}
% Objetos experimentais, população e amostra (histórias do CAPJu).

O experimento será realizado no contexto de dois projetos de aplicações web desenvolvidos com métodos ágeis na disciplina Engenharia de Produto de Software da Universidade de Brasília, em diferentes semestres: o MeasureSoftGram, desenvolvido no semestre 2026/2 \cite{measuresoftgram2026doc}, e o CAPJu, desenvolvido no semestre 2023/1 \cite{capju2023doc}. Para ambos os projetos, serão utilizados os artefatos disponíveis em seus respectivos repositórios, incluindo documentações, histórias de usuário e protótipos.

Além do repositório original do CAPJu \cite{capju2023doc}, será utilizado outro repositório, que contém cenários de aceitação escritos em Gherkin por Gustavo França Boa Sorte e Leandro de Almeida Oliveira, que foram elaborados a partir das histórias de usuário do projeto, disponíveis no repositório do Trabalho de Conclusão de Curso dos autores citados \cite{capju2026cenarios} . Assim, para o CAPJu, serão utilizados os artefatos do repositório original e os cenários Gherkin oriundos das histórias de usuário.

A escolha destes projetos é feita porque existem histórias de usuário com critérios de aceitação descritos em BDD e cuja validação pelos usuários é conhecida. Isso é relevante porque podemos utilizar os critérios de aceitação definidos e validados como referência para avaliar os cenários Gherkin elaborados. Como não existem muitos dados públicos com essas características, foram escolhidos projetos cujos artefatos estão disponíveis e cuja origem e validação podem ser rastreadas. Portanto, a seleção é uma amostragem por conveniência, cujas limitações são discutidas na seção~\ref{sec:validade} de ameaças à validade deste trabalho.

\textbf{Unidades experimentais:} histórias de usuário dos dois projetos, submetidas aos tratamentos com e sem RAG para a geração de critérios de aceitação em formato Gherkin.

\textbf{População:} histórias de usuário de aplicações web desenvolvidas com métodos ágeis que possuam documentação de domínio e critérios de aceitação definidos e validados.

\textbf{Amostra:} 50 histórias de usuário, sendo 27 do MeasureSoftGram e 23 do CAPJu.

No MeasureSoftGram, os critérios de aceitação estão descritos no formato Given-When-Then, mas não estão organizados em arquivos Gherkin. Então, esses critérios serão transcritos para a sintaxe Gherkin, mantendo o conteúdo original e fazendo apenas as adaptações estruturais necessárias. A conformidade sintática dos arquivos resultantes será verificada utilizando o framework Python Behave. Os cenários elaborados pelos modelos, tanto no tratamento com RAG quanto no tratamento sem RAG, também serão verificados, de modo que possíveis problemas de sintaxe sejam identificados separadamente da avaliação de acurácia. Esta última será feita utilizando o conteúdo dos critérios elaborados, enquanto a verificação com o Behave avaliará sua conformidade com a sintaxe Gherkin.

Os critérios de aceitação utilizados como gabarito em cada projeto serão decompostos em objetivos de teste atômicos, seguindo a abordagem apresentada por \cite{wang2025}. Estes objetivos serão utilizados como unidades de referência para a avaliação da cobertura dos critérios elaborados. As histórias de usuário serão submetidas a prompts padronizados em dois tratamentos: com RAG e sem RAG. Para cada história, os resultados dos dois tratamentos serão avaliados em relação ao mesmo gabarito, permitindo a comparação da acurácia nas duas condições experimentais.

\subsection{Variáveis do experimento}
% {aqui vão as variáveis independentes (ex.: LLM, técnica de prompting), dependentes (ex.: acurácia dos critérios) e as hipóteses nula e alternativa, porque definem o que o experimento testa}

\subsubsection*{Fator experimental}
% Estratégia de contexto, com dois níveis (sem RAG e com RAG).
O experimento apresenta apenas um fator: \textbf{Uso do RAG para a geração de critérios de aceitação no Gherkin}. O fator possui dois níveis (tratamentos) relativos a duas configurações do mesmo pipeline de geração.

\begin{enumerate} 
    \item \textbf{Sem RAG}: o modelo é fornecido com o prompt padronizado e com a história de usuário; 
    \item \textbf{Com RAG}: o modelo é fornecido com o mesmo prompt e com a mesma história, acompanhados das informações de domínio recuperadas da base de conhecimento do projeto. 
\end{enumerate}

O \textit{pipeline} será de implementação própria, baseado na arquitetura RAGcceptance\~M2RE \cite{wang2025}, e será segmentado em duas partes: recuperação e geração. O prompt será elaborado a partir da técnica de \textit{prompting} que apresentou melhor desempenho em \citeonline{wang2025} (a estrutura \textit{Urial}), que consiste em um template fixo com campos para a história de usuário e para o contexto do projeto.

A parte da recuperação apresentará dois módulos: a recuperação textual (T-RAG), que irá fornecer os cinco trechos mais relevantes da documentação por ICRALM, e a recuperação visual (V-RAG), que irá extrair as capturas de tela da interface ligadas por DSE \cite{wang2025}. No tratamento com RAG, o campo de contexto do template será preenchido com as informações e com as imagens recuperadas. No tratamento sem RAG, este campo ficará vazio, enquanto os outros elementos do template não serão alterados.

Os dois tratamentos irão utilizar o mesmo modelo de linguagem para a geração (e.g., Gemini) \cite{wang2025}, o mesmo template de prompt e os mesmos parâmetros de geração com temperatura zero. Dessa forma, a única diferença entre os tratamentos será o conteúdo do campo de contexto.

O tratamento sem RAG será idêntico ao cenário w/o RAG de \citeonline{wang2025}, onde o template é mantido, mas o contexto de fundo é removido.

\subsubsection*{Variável independente}
% TODO

O tratamento aplicado é a variável independente, que se relaciona à estratégia de geração de critérios de aceitação em Gherkin usada (\textbf{Sem RAG} ou \textbf{Com RAG}), segundo os dois níveis estipulados no fator experimental.

\subsubsection*{Variáveis dependentes}
% Acurácia (Acc@Hit e Acc@Cor); similaridade como medida complementar.

A variável dependente primária é a precisão dos critérios de aceitação gerados em Gherkin, que são medidos com base na cobertura dos objetivos atômicos do modelo de referência (\textit{LLM-as-a-Judge}). Para cada objetivo atômico ($T_i$) extraído do modelo de uma história de usuário, o grupo de LLMs juízes dá o nível de cobertura em \textit{Full Coverage}, \textit{Partial Coverage} ou \textit{Not Covered} \cite{wang2025}. Os níveis \textit{Full Coverage} e \textit{Not Covered} só são dados quando há unanimidade entre os juízes; se não houver unanimidade, o objetivo é classificado como \textit{Partial Coverage}. Com base nessa classificação, são definidos duas métricas de precisão principais, calculadas para cada história de usuário:
\begin{itemize} 
    \item \textbf{Taxa de Cobertura de Atendimento (\textit{Acc@Hit})}: percentual de objetivos atômicos da história de usuário que atingiram pelo menos cobertura parcial ($Hit = 1$); 
    \item \textbf{Taxa de Cobertura de Exatidão (\textit{Acc@Cor})}: percentual de objetivos atômicos da história de usuário que atingiram cobertura total ($Correctness = 1$). 
\end{itemize} 
Com base nestas métricas, serão obtidas as seguintes estatísticas por tratamento: 
\begin{itemize} 
    \item \textbf{Média} e \textbf{mediana} das taxas de cobertura (\textit{Acc@Hit} e \textit{Acc@Cor}) para cada história de usuário; 
    \item \textbf{Variância} e \textbf{desvio padrão} das distribuições de precisão entre as histórias de usuário em cada tratamento. 
\end{itemize}

A similaridade semântica entre os critérios gerados e os critérios do modelo será calculada como medida extra e apresentada de maneira descritiva, de acordo com a Seção~\ref{sec:analise-complementar}.

\subsubsection*{Variáveis controladas}
% LLM e versão, temperatura, prompt, corpus, top-k.
As variáveis controladas são os fatores que se mantiveram constantes entre os dois tratamentos (\textbf{Sem RAG} e \textbf{Com RAG}) para que o efeito da incorporação das informações de domínio fosse isolado e não houvesse vieses no experimento.
Estas variáveis são:
\begin{itemize}
    \item \textbf{Modelo de LLM gerador e versão}: o mesmo modelo de linguagem gerador (por exemplo, Gemini 2 Flash) é empregado para as duas configurações de uma história de usuário.
    \item \textbf{Parâmetros de geração}: os hiperparâmetros de inferência são fixados em ambos os tratamentos. Mais precisamente, a temperatura é igual a zero ($temperatura = 0$) e $top\_p = 0.1$. Com isso, a aleatoriedade e variabilidade das respostas são reduzidas, aumentando a reprodutibilidade dos resultados (mesmo não garantindo que as saídas serão as mesmas entre as execuções).
    \item \textbf{Estrutura do template de prompt}: a mesma estrutura de prompt padronizada (adaptada da técnica \textit{Urial}) é utilizada. O papel da IA, as instruções do formato Gherkin e o espaço para a história de usuário são idênticos.
    \item \textbf{Corpus da base de conhecimento}: a mesma base de documentos textuais e visuais do projeto é usada para indexação e busca.
    \item \textbf{Quantidade de trechos recuperados}: cinco trechos da documentação e telas são retornados pelo módulo de recuperação no tratamento Com RAG.
    \item \textbf{Comitê e avaliação (\textit{LLM-as-a-Judge})}: o mesmo painel de LLMs juízas (GPT-4o, Claude 3.5 Sonnet e Gemini 2 Flash) é usado. Ele é submetido ao mesmo algoritmo de julgamento ponto a ponto e à exigência de consenso unânime para classificar as saídas de ambos os tratamentos.
\end{itemize}

\subsection{Formulação das hipóteses}
% H0 e H1, nível de significância.

As hipóteses para a Questão Específica 1 da seção~\ref{sec:questao-pesquisa} serão modeladas como uma diferença pareada entre os tratamentos, de forma a quantificar, estatisticamente, o impacto do uso de RAG na geração de critérios de aceitação em formato Gherkin. Para cada história de usuário $j$, temos:
\[
d_j^{\mathrm{Hit}} = \mathrm{Acc@Hit}_{j,\mathrm{Com\ RAG}} - \mathrm{Acc@Hit}_{j,\mathrm{Sem\ RAG}}
\quad \mathrm{e} \quad
d_j^{\mathrm{Cor}} = \mathrm{Acc@Cor}_{j,\mathrm{Com\ RAG}} - \mathrm{Acc@Cor}_{j,\mathrm{Sem\ RAG}}
\]
como as diferenças nas taxas obtidas com e sem RAG e $\mu_d^{\mathrm{Hit}}$ e $\mu_d^{\mathrm{Cor}}$ as respectivas médias populacionais dessas diferenças \cite{wohlin2000}.

\subsubsection*{1. Hipóteses para a Taxa de Cobertura de Atendimento (\textit{Acc@Hit})}

\begin{itemize}
    \item \textbf{$H_{0,\mathrm{Hit}}$ (Hipótese Nula)}: a geração \textbf{Com RAG} e a geração \textbf{Sem RAG} não diferem quanto à taxa de cobertura de atendimento (\textit{Acc@Hit}).

    \[
    H_{0,\mathrm{Hit}}: \mu_d^{\mathrm{Hit}} = 0
    \]

    \item \textbf{$H_{1,\mathrm{Hit}}$ (Hipótese Alternativa)}: a geração \textbf{Com RAG} e a geração \textbf{Sem RAG} diferem quanto à taxa de cobertura de atendimento (\textit{Acc@Hit}).

    \[
    H_{1,\mathrm{Hit}}: \mu_d^{\mathrm{Hit}} \neq 0
    \]
\end{itemize}

\subsubsection*{2. Hipóteses para a Taxa de Cobertura de Exatidão (\textit{Acc@Cor})}

\begin{itemize}
    \item \textbf{$H_{0,\mathrm{Cor}}$ (Hipótese Nula)}: a geração \textbf{Com RAG} e a geração \textbf{Sem RAG} não diferem quanto à taxa de cobertura de exatidão (\textit{Acc@Cor}).

    \[
    H_{0,\mathrm{Cor}}: \mu_d^{\mathrm{Cor}} = 0
    \]

    \item \textbf{$H_{1,\mathrm{Cor}}$ (Hipótese Alternativa)}: a geração \textbf{Com RAG} e a geração \textbf{Sem RAG} diferem quanto à taxa de cobertura de exatidão (\textit{Acc@Cor}).

    \[
    H_{1,\mathrm{Cor}}: \mu_d^{\mathrm{Cor}} \neq 0
    \]
\end{itemize}

Se a hipótese nula for rejeitada, o sinal da diferença média indicará qual dos tratamentos obteve maior cobertura. Quando a normalidade das diferenças for rejeitada e o teste de Wilcoxon for utilizado, a hipótese nula será avaliada de forma semelhante, verificando se a distribuição das diferenças é zero.

% TODO: decidir se haverá correção para múltiplos testes (duas métricas primárias: Acc@Hit e Acc@Cor), por exemplo Holm-Bonferroni.

\subsubsection*{Nível de Significância e Testes Estatísticos}

O nível de significância é escolhido em $\alpha = 0,05$ (nível de confiança de 95\%), padrão em Engenharia de Software Experimental \cite{wohlin2000}, para as hipóteses. Se o valor-p (\textit{p-value}) que é obtido através dos testes de hipótese for menor que $\alpha$ ($p < 0,05$), então a hipótese nula que está associada será rejeitada.

\subsection{Desenho do experimento}
% {aqui vão os tratamentos e a forma de alocação, porque determina como os dados serão comparados}
% Um fator, dois tratamentos, comparação pareada.
Com base na tipologia de \citeonline{wohlin2000}, o experimento descrito aqui se enquadra na categoria de \textbf{experimento orientado à tecnologia} (\textit{technology-oriented experiment}). Nesse tipo de pesquisa não há a participação de seres humanos, e o intuito é comparar experimentalmente técnicas, ferramentas ou algoritmos aplicados a artefatos de \textit{software}. No nosso caso, a geração de critérios de aceitação em formato Gherkin a partir de histórias de usuário.

O desenho experimental escolhido é o \textbf{\textit{paired comparison design}} (desenho de comparação pareada) \cite{wohlin2000}. Essa abordagem é indicada para experimentos com um fator e dois tratamentos. Cada unidade experimental (\textit{User Story}) recebe os dois tratamentos (\textbf{Sem RAG} e \textbf{Com RAG}), e a comparação é feita sobre a diferença entre as medidas pareadas. A ordem de aplicação dos tratamentos é escolhida aleatoriamente para cada história. Como as chamadas aos modelos são independentes e não há seres humanos que possam aprender com a primeira execução, a aleatorização não serve para controlar efeitos de aprendizado. Serve para evitar que um dos tratamentos seja executado sistematicamente antes do outro, por exemplo, caso haja mudanças de disponibilidade ou de versão dos serviços dos modelos durante a coleta.

A Tabela~\ref{tab:pareado-exemplo} dá um exemplo de ordem aleatória de execução dos tratamentos sobre uma amostra de histórias de usuário.

\begin{table}[H]
\centering
\caption{Exemplo ilustrativo da ordem aleatória de aplicação dos tratamentos às histórias de usuário (\textit{paired comparison design})}
\label{tab:pareado-exemplo}
\begin{tabular}{c|c|c}
\hline
\textbf{Unidade experimental (História de Usuário)} & \textbf{Trat. 1 (Sem RAG)} & \textbf{Trat. 2 (Com RAG)} \\
\hline
1 & 2 & 1 \\
2 & 1 & 2 \\
3 & 2 & 1 \\
4 & 1 & 2 \\
5 & 2 & 1 \\
\hline
\end{tabular}

\vspace{0.3em}
\small{Legenda: Os valores nas células indicam a ordem de execução do tratamento para a história correspondente. Tratamentos: 1 = Sem RAG (\textit{baseline}), 2 = Com RAG (tratamento proposto).}
\end{table}

\subsection{Teste de normalidade}

Antes da comparação estatística entre os tratamentos (\textbf{Sem RAG} e \textbf{Com RAG}) ser feita, o pressuposto de normalidade das diferenças pareadas entre os tratamentos ($d_j^{\mathrm{Hit}}$ e $d_j^{\mathrm{Cor}}$), calculadas para cada história de usuário, será verificado por meio do teste de Shapiro-Wilk. Esse teste é apropriado para amostras de pequeno a médio tamanho, como o conjunto de 50 pares do experimento em análise.

Na Engenharia de Software Experimental, é comum que dados sobre qualidade e acurácia de requisitos apresentem assimetria ou não atendam aos pressupostos necessários para testes paramétricos \cite{wohlin2000}. No caso de rejeição da hipótese de normalidade ($p < 0{,}05$), o estudo seguirá a abordagem de testes não paramétricos para amostras pareadas, utilizando o \textbf{teste de postos sinalizados de Wilcoxon} (\textit{Wilcoxon signed-rank test}), com a medida do tamanho do efeito (\textit{effect size}) pelo $r$ de Rosenthal. No caso contrário, será empregado o \textbf{teste $t$ pareado}, com o $d$ de Cohen.

\subsection{Teste de hipóteses}
% {aqui vão as métricas e os métodos de análise (estatística descritiva, testes), porque a escolha deve vir antes da coleta}
% Teste t pareado ou Wilcoxon.

O estudo estatístico será efetuado a nível da história de usuário (\textit{case-level}), uma vez que a unidade experimental é a história de usuário sobre a qual foram feitas as hipóteses e sobre a qual ambos os tratamentos (\textbf{Sem RAG} e \textbf{Com RAG}) são realizados. Para cada história, as diferenças pareadas $d_j^{\mathrm{Hit}}$ e $d_j^{\mathrm{Cor}}$ serão calculadas. No caso em que não há rejeição da hipótese de normalidade destas diferenças, o \textbf{teste $t$ pareado} será utilizado. Caso contrário, o \textbf{teste de postos sinalizados de Wilcoxon} (\textit{Wilcoxon signed-rank test}) \cite{wohlin2000} será utilizado. Ambos os testes são adequados para a comparação de duas medidas pareadas sobre as mesmas unidades experimentais.

As descobertas no nível do objetivo atômico (\textit{point-level}) serão relatadas de maneira descritiva, sendo a contagem de objetivos cobertos em cada tratamento. Não serão usadas para testar as hipóteses, pois objetivos atômicos derivados da mesma história não são observações independentes.

\subsection{Instrumentação}
% {aqui vão as ferramentas, os prompts, o gabarito de referência e os scripts, porque permite reproduzir o estudo}

A experimentação possui como instrumentos de trabalho uma série de recursos em termos de \textit{software}, templates de \textit{prompt}, gabarito de referência, scripts de automação, artefatos, entre outros, para permitir a execução dos tratamentos e a coleta de dados de modo a tornar o estudo inteiramente reprodutível. A organização dos instrumentos segue a estrutura abaixo:

\begin{itemize}
    \item \textbf{Gabarito de Referência e Base de Conhecimento}: É composto por um conjunto de histórias de usuário (\textit{User Stories}) reais do projeto, com seus respectivos critérios de aceitação em formato Gherkin (\textit{ground truth}), que no caso do MeasureSoftGram são transcritos dos critérios originais descritos em \textit{Given-When-Then}. Para a avaliação, os cenários em Gherkin do gabarito serão divididos em uma lista de objetivos de teste atômicos ($T_i$), mantendo a semântica condicional (\textit{Given-When-Then}). A base de conhecimento é composta por documentos textuais e capturas de tela da interface do usuário (\textit{UI}) indexados para recuperação de domínio;

    \item \textbf{Ferramentas de Software e Modelos}: Uma infraestrutura construída na linguagem Python (versão 3.12) para gerenciar o pipeline. Para a geração, usa-se a API da LLM geradora (\textit{e.g.}, Gemini 2 Flash). Para a avaliação (\textit{LLM-as-a-Judge}), usa-se as APIs dos três modelos que compõem o comitê de juízes (GPT-4o, Claude 3.5 Sonnet e Gemini 2 Flash) \cite{wang2025};

    \item \textbf{Templates de Prompt}: O template de prompt é padronizado e fixo, baseado na estrutura \textit{Urial} \cite{wang2025}. Ele define o papel da IA, as regras de sintaxe Gherkin e tem lacunas para incluir a história de usuário (\texttt{\{INPUT\_USER\_STORY\}}) e o contexto recuperado (\texttt{\{INPUT\_BACKGROUND\}}). No tratamento Sem RAG, o espaço para contexto fica em branco, enquanto no tratamento Com RAG ele é preenchido com os cinco trechos recuperados;

    \item \textbf{Scripts de Automação e Recuperação}: Scripts para extrair e vetorizar artefatos da base de conhecimento, realizar a busca por similaridade e fazer as chamadas à LLM geradora de maneira pareada para cada história de usuário, respeitando a ordem aleatória do desenho de comparação pareada;

    \item \textbf{Scripts de Avaliação e Análise Estatística}: Scripts em Python para rodar o algoritmo de julgamento \textit{LLM-as-a-Judge}, enviando para o painel de juízes cada objetivo atômico $T_i$ e a resposta gerada, coletando o consenso de cobertura (\textit{Full}, \textit{Partial} ou \textit{Not Covered}), computando as métricas de \textit{Acc@Hit}, \textit{Acc@Cor} e similaridade semântica, e conduzindo o teste de normalidade e os testes de hipóteses (Shapiro-Wilk, teste $t$ pareado e Wilcoxon);

    \item \textbf{Behave}: Framework Python para BDD utilizado para verificar a sintaxe correta dos arquivos Gherkin, tanto os critérios do MeasureSoftGram convertidos para esta sintaxe quanto os cenários gerados pelos modelos nos dois tratamentos. Os resultados dessa verificação são armazenados de forma independente da avaliação de acurácia, realizada com base no conteúdo dos critérios.
\end{itemize}

\subsection{Procedimento de aplicação dos instrumentos}
% {aqui vão os passos de execução em ordem, porque a etapa do TCC2 vai seguir esse roteiro}

O método de utilização dos instrumentos é o plano de execução passo a passo do experimento que será feito no Trabalho de Conclusão de Curso 2. O processo é organizado em sete fases seguidas, permitindo acompanhar os dados e manter as variáveis do experimento sob controle:

\begin{enumerate}
    \item \textbf{Preparação da Base de Conhecimento e Indexação (Pré-processamento)}:
    \begin{itemize}
        \item Pesquisa e organização das histórias de usuário (\textit{User Stories}) e de suas respectivas informações (textos e imagens da tela do usuário);
        \item Indexação de vetores da base de conhecimento usando o modelo ICRALM para obter trechos de texto (\textit{T-RAG}) e do modelo DSE para obter imagens da tela do usuário (\textit{V-RAG}) \cite{wang2025}.
    \end{itemize}

    \item \textbf{Transformação do Gabarito em Gherkin}:
    \begin{itemize}
        \item Conversão dos critérios de aceitação do MeasureSoftGram (dados em \textit{Given-When-Then}) para a linguagem Gherkin, sem alterar o conteúdo e fazendo com que os arquivos gerados tenham a linguagem do \textit{Behave};
        \item Separação manual/organizada dos critérios de aceitação do gabarito de cada história de usuário em uma lista de objetivos de teste menores ($T_i$), mantendo a regra de negócio intacta (\textit{Given-When-Then});
        \item Salvamento dos objetivos menores em banco de dados organizado para usar na hora de avaliar.
    \end{itemize}

    \item \textbf{Sorteio da Ordem dos Tratamentos}:
    \begin{itemize}
        \item Uso do algoritmo de sorteio para decidir em qual ordem serão feitos os dois tratamentos (\textbf{Sem RAG} e \textbf{Com RAG}) para cada história de usuário, conforme o jeito de comparar em pares (\textit{paired comparison design}) mostrado na Tabela~\ref{tab:pareado-exemplo};
        \item Salvamento da ordem de execução em repositório de controle para poder seguir a execução.
    \end{itemize}

    \item \textbf{Execução dos Tratamentos e Geração dos Gherkins}:
    \begin{itemize}
        \item \textit{Tratamento Com RAG}: Para a história de usuário em execução, o módulo de recuperação extrai os cinco fragmentos textuais e visuais mais significativos da base de conhecimento. O contexto obtido é colocado no \textit{template} de \textit{prompt} (com base no modelo \textit{Urial}), e a solicitação é feita à LLM geradora (\textit{e.g.}, Gemini 2 Flash) com a temperatura em 0;
        \item \textit{Tratamento Sem RAG}: O mesmo \textit{template} de \textit{prompt} e história de usuário são enviados à mesma LLM geradora com os mesmos hiperparâmetros, mas o campo de contexto está em branco;
        \item Gravação das respostas com os cenários em Gherkin criados para ambos os tratamentos em arquivos JSON marcados pelo ID da história e tratamento.
    \end{itemize}

    \item \textbf{Análise Sintática dos Cenários Criados (\textit{Behave})}:
    \begin{itemize}
        \item Conferência da sintaxe Gherkin dos cenários criados em cada tratamento, usando o \textit{Behave};
        \item Registro do resultado da conferência separado da avaliação de acurácia.
    \end{itemize}

    \item \textbf{Avaliação Automática por meio do Comitê de LLMs Juízas (\textit{LLM-as-a-Judge})}:
    \begin{itemize}
        \item Para cada objetivo atômico $T_i$ de cada história de usuário, os \textit{scripts} enviam o objetivo e a resposta completa criada pelo tratamento ao painel de juízes (GPT-4o, Claude 3.5 Sonnet e Gemini 2 Flash);
        \item Obtenção dos pareceres individuais de cada juiz e uso do critério de consenso: as classificações \textit{Full Coverage} e \textit{Not Covered} são dadas somente quando todos os juízes concordam, e se houver discordância, o objetivo é classificado como \textit{Partial Coverage};
        \item Associação do resultado para as variáveis binárias de atendimento ($Hit$ e $Correctness$) por objetivo atômico.
    \end{itemize}

    \item \textbf{Cálculo das Métricas e Análise Estatística}:
    \begin{itemize}
        \item Agregação dos dados para calcular as taxas \textit{Acc@Hit} e \textit{Acc@Cor} em nível de caso/história (\textit{Case-level}) e similaridade semântica, assim como a contagem descritiva de objetivos atingidos em nível de ponto (\textit{Point-level});
        \item Aplicação do teste de normalidade de Shapiro-Wilk nas diferenças pareadas entre os tratamentos;
        \item Realização do teste de hipóteses pareado em nível de caso (teste $t$ pareado ou teste de Wilcoxon, de acordo com a normalidade das diferenças) ao nível de significância de $\alpha = 0{,}05$, acompanhado pelo cálculo do tamanho do efeito ($d$ de Cohen para o teste $t$ pareado ou $r$ de Rosenthal para o teste de Wilcoxon).
    \end{itemize}
\end{enumerate}

\subsection{Validade do experimento}
\label{sec:validade}


Seguindo as diretrizes dos protocolos de Engenharia de Software Experimental, o planejamento do experimento deve reconhecer ativamente as ameaças potenciais à validade dos resultados e estabelecer estratégias de mitigação para cada classe.

\subsubsection*{Validade de construto}

A validade de construto examina se as métricas escolhidas representam adequadamente o conceito abstrato de acurácia dos critérios de aceitação em Gherkin.

\begin{itemize}
    \item \textbf{Ameaça}: As métricas estatísticas de superfície mais comuns (como ROUGE, BLEU e Levenshtein) não podem capturar o significado de regras de negócio em formato Gherkin. Por outro lado, o uso de juízes baseados em LLMs pode introduzir viés de alinhamento, sensibilidade ao \textit{prompt} ou inconsistência de avaliação.

    \item \textbf{Mitigação}: Uma abordagem complementar de métricas é utilizada. O processo de avaliação \textit{LLM-as-a-Judge} é conduzido sobre objetivos de teste atômicos ($T_i$) obtidos a partir da desconstrução prévia do gabarito. Para mitigar os vieses individuais dos modelos, é empregado um comitê diversificado de três LLMs juízas (GPT-4o, Claude 3.5 Sonnet e Gemini 2 Flash) sujeitas a um critério de \textbf{consenso unânime} para as classificações \textit{Full Coverage} e \textit{Not Covered}, além do cálculo da similaridade semântica por \textit{embeddings}, que é apresentado de forma complementar. Como a falta de unanimidade resulta em \textit{Partial Coverage}, a métrica \textit{Acc@Hit} tende a ser mais permissiva que a \textit{Acc@Cor}, o que é levado em conta na interpretação dos resultados.
    % TODO: decidir se haverá validação do comitê de juízes por avaliação humana de uma subamostra; se não, declarar a ausência como limitação.

    \item \textbf{Ameaça}: Viés de preferência por texto próprio, já que o modelo gerador pode coincidir com um dos modelos do comitê de juízes (como o Gemini).
    % TODO: definir o modelo gerador (hoje aparece como e.g.).

    \item \textbf{Mitigação}: O comitê é formado por três juízes, nenhum dos quais decide sozinho. Além disso, o mesmo gerador é utilizado nos dois tratamentos, de modo que um eventual viés tende a afetar de forma semelhante as gerações Com RAG e Sem RAG, comprometendo mais o nível absoluto de acurácia do que a diferença pareada entre os tratamentos, que é o foco do experimento.

    \item \textbf{Ameaça}: \textit{Acc@Hit} e \textit{Acc@Cor} avaliam a cobertura dos objetivos atômicos do gabarito. Os critérios gerados não mapeados a nenhum desses objetivos, incluindo critérios espúrios criados pelo modelo, não mudam os valores das métricas. Da mesma forma, os critérios válidos não existentes no gabarito não são considerados acertos, pois o gabarito não possui todos os possíveis critérios válidos.

    \item \textbf{Mitigação}: Esta limitação é assumida e levada em conta na interpretação dos resultados, que são relativos à cobertura do gabarito, não à inexistência de critérios espúrios. A análise qualitativa complementar (Seção~\ref{sec:analise-complementar}) analisa os tipos de falha, incluindo alucinações.
    % TODO: considerar amostrar a análise qualitativa de alucinações em todas as histórias, não só nas de baixa pontuação.

    \item \textbf{Ameaça}: Alguns do gabarito são criados pelos autores, visto que os critérios do MeasureSoftGram serão transformados para a sintaxe Gherkin, enquanto os cenários do CAPJu foram criados por outros, o que pode trazer vieses de redação ou omissões.

    \item \textbf{Mitigação}: A transformação mantém o conteúdo original dos critérios, somente as adaptações estruturais necessárias são feitas, e a conformidade sintática é verificada com o \textit{Behave}.
\end{itemize}

\subsubsection*{Validade interna}

A validade interna é o grau de confiabilidade com que podemos dizer que os efeitos observados nas variáveis dependentes são devidos unicamente à manipulação da variável independente (\textbf{Uso de RAG}) e não a outros fatores confundentes.

\begin{itemize}
    \item \textbf{Ameaça}: O não-determinismo intrínseco à geração por modelos de linguagem de grande porte, variações de contexto e possíveis mudanças na disponibilidade ou na versão dos serviços durante a coleta de dados.

    \item \textbf{Mitigação}: Nós fixamos os hiperparâmetros de inferência com temperatura zerada ($\mathrm{temperatura} = 0$) e $top\_p = 0{,}1$ em todas as execuções do gerador e dos juízes, reduzindo assim a variabilidade das respostas, embora não garantindo saídas idênticas entre execuções. Usamos um desenho de comparação pareada, no qual cada história de usuário recebe ambos os tratamentos usando o mesmo modelo gerador e a mesma estrutura de \textit{prompt} (\textit{Urial}), e a ordem de aplicação é aleatorizada para que um dos tratamentos não seja executado sistematicamente antes do outro.
\end{itemize}

\subsubsection*{Validade externa}

A validade externa refere-se à possibilidade de generalizar os resultados para outros contextos, modelos e sistemas de software.

\begin{itemize}
    \item \textbf{Ameaça}: Nossa amostra é de conveniência e consiste em 50 histórias de usuário ($N = 50$), provenientes de apenas dois projetos acadêmicos de uma mesma disciplina, o que pode restringir a generalização dos resultados para projetos de grande porte, sistemas em produção ou com padrões de documentação distintos.

    \item \textbf{Mitigação}: Os artefatos dos projetos selecionados estão consolidados, já que são projetos acadêmicos terminados, e sua origem e validação podem ser rastreadas. Ainda assim, a generalização dos resultados continua restrita a esse contexto. A replicação em outros domínios e projetos é facilitada pelo fato de disponibilizarmos publicamente o \textit{pipeline} de automação, os \textit{templates} de \textit{prompt} e os \textit{scripts} de avaliação em um repositório de artefatos.
\end{itemize}

\subsubsection*{Validade de conclusão}

A validade de conclusão refere-se à relação estatística entre a variável independente e os resultados, assegurando a confiabilidade dos testes empregados e a capacidade de realizar inferências.

\begin{itemize}
    \item \textbf{Ameaça}: Perigo de perda de poder estatístico (\textit{statistical power}) ou de infringimento dos pressupostos distributivos devido à restrição do tamanho amostral em 50 histórias de usuário ($N = 50$), o que poderia causar erros de Tipo I (falsos positivos) ou Tipo II (falsos negativos).

    \item \textbf{Mitigação}: O desenho de comparação pareada é o principal atenuante, pois cada uma das 50 histórias de usuário é submetida aos dois tratamentos (\textbf{Sem RAG} e \textbf{Com RAG}), gerando 50 pares de comparação direta (100 execuções de geração no total) e eliminando a variabilidade intrínseca a cada história. A inferência é realizada ao nível da história, uma vez que os objetivos atômicos de uma mesma história não são observações independentes. Para testar a normalidade das diferenças pareadas, utiliza-se o \textbf{teste de Shapiro-Wilk}, adequado para amostras de pequeno a médio porte. Se a normalidade for negada, utiliza-se o teste de Wilcoxon; caso contrário, o teste $t$ pareado, sob o nível de significância de $\alpha = 0{,}05$, acompanhados pelo cálculo do tamanho do efeito (\textit{effect size}, através de $r$ de Rosenthal para o teste de Wilcoxon e $d$ de Cohen para o teste $t$ pareado) para medir a importância prática das diferenças observadas.
\end{itemize}

\subsection{Análise complementar}
\label{sec:analise-complementar}
% Medidas complementares (ex.: similaridade de cosseno), de caráter exploratório.

Além da verificação das hipóteses primárias que dizem respeito à taxa de cobertura (\textit{Acc@Hit} e \textit{Acc@Cor}), será feito um estudo descritivo e exploratório, sem testes de hipóteses, com o objetivo de dar uma visão geral da estrutura dos critérios gerados e dos tipos de erros mais comuns em cada tratamento:

\begin{enumerate}
    \item \textbf{Similaridade Semântica Vetorial}: similaridade de cosseno entre os \textit{embeddings} dos critérios em Gherkin gerados e os critérios do gabarito da história de usuário correspondente \cite{wang2025}. Esta medida será apresentada descritivamente (média e mediana por tratamento), sem hipótese estatística associada, pois a análise estatística de medidas automáticas não considera o contexto dos critérios \cite{wang2025};

    \item \textbf{Métricas Estatísticas e Lexicais Secundárias}: Para avaliar a sobreposição sintática e a similaridade estrutural entre os critérios em Gherkin gerados e o gabarito de referência, serão calculadas as seguintes métricas automáticas quantitativas \cite{wang2025}:
    \begin{itemize}
        \item \textbf{ROUGE (ROUGE-1, ROUGE-2 e ROUGE-L)}: para medir a sobreposição de unigramas, bigramas e da maior sequência comum de palavras (\textit{Longest Common Sequence});
        \item \textbf{BLEU}: para calcular a precisão de n-gramas com penalização por brevidade (\textit{brevity penalty});
        \item \textbf{Distância de Levenshtein}: para medir o menor número de edições de caracteres necessárias para alinhar os cenários gerados ao gabarito.
    \end{itemize}

    \item \textbf{Análise Qualitativa de Modos de Falha e Alucinações}: Inspeção qualitativa da amostra de cenários que receberam baixas pontuações pelos juízes (\textit{Not Covered} ou \textit{Partial Coverage}), agrupando as falhas em três categorias principais:
    \begin{itemize}
        \item \textit{Alucinação de Regras}: inclusão de pré-condições ou asserções falsas não baseadas na história ou na documentação;
        \item \textit{Omissão de Fluxos de Exceção}: falta de critérios que cobram caminhos negativos e restrições de interface;
        \item \textit{Redundância e Verbosidade}: criação de critérios duplicados ou paráfrases de uma mesma condição.
    \end{itemize}
\end{enumerate}